\documentclass[
	letterpaper, % Use US letter paper size
]{jdf}

\usepackage{todonotes}
\addbibresource{references.bib}

\author{Your Name}
\email{youremail@gatech.edu}
\title{KBAI Assignment:\\Getting to Know Jill Watson}

\begin{document}

\maketitle

% =============================================================================
% The abstract is OPTIONAL. Include one only if your paper benefits from a
% high-level summary. If you do not need it, delete the abstract environment
% below. If you keep it, replace the placeholder text with your own summary.
% =============================================================================
\begin{abstract}
	(Optional) Replace this text with a brief summary of your findings from
	interacting with Jill Watson across all three parts of this assignment.
	Delete this entire abstract block if you choose not to include one.
\end{abstract}

% =============================================================================
% JDF FORMATTING REFERENCE (this comment block is for your reference only)
% =============================================================================
% TYPOGRAPHY: Palatino, 11 pt body, 17 pt line spacing, 8.5 pt after para.
% TITLE: 17 pt, centered, up to 3 lines.
% HEADINGS: 11 pt bold. H1 ALL CAPS. H2 bold roman. H3 bold italic.
%   H4 run-in sidehead (bold italic + em dash). Sentence case except H1.
% PAGE LAYOUT: US Letter, 1" top margin, 1.5" bottom and sides.
% FIGURES: Centered, caption below, 8.5 pt caption, referenced in text.
% TABLES: Caption above ("figures at the foot, tables at the top").
% QUOTES: <3 lines in-line; 3+ lines as blockquote (0.5" extra margins).
% LISTS: Indented 0.5" from left, bullet/number hanging 0.25".
% CITATIONS: APA preferred. \citep{}, \citeauthor{}, \citeyear{}.
% FOOTNOTES: 8.5 pt, 14 pt line spacing. In-line citations preferred.
% REFERENCES: Numbered, alphabetical, only cited works. Doesn't count
%   against page limit.
% =============================================================================


% #############################################################################
%                   PART 1: CONCEPT EXPLORATIONS
% #############################################################################
\section{Concept explorations}

% ~~~~~~~~~~~~~~~~~~~~~~~~~~~~~~~~~~~~~~~~~~~~~~~~~~~~~~~~~~~~~~~~~~~~~~~~~~~~~
% INSTRUCTIONS (for your reference — not rendered in the PDF):
%
% Complete three targeted interactions with Jill Watson (~2.5 hours total).
% For each activity you are responsible for formulating your OWN prompts.
% What you are given is a scenario (the context) and a conceptual target
% (what the conversation must illuminate). How you get there is the exercise.
%
% Use a FRESH JW session for each activity to prevent cross-contamination.
% Design your scenarios BEFORE opening JW — asking JW to generate your
% scenario undermines the purpose of the activity.
%
% Each response should be 2–3 paragraphs.
%
% Include one JW chat transcript per activity in the Appendix (3 total).
% ~~~~~~~~~~~~~~~~~~~~~~~~~~~~~~~~~~~~~~~~~~~~~~~~~~~~~~~~~~~~~~~~~~~~~~~~~~~~~

\subsection{Activity A: The four schools and grand challenges (Chapter 1)}

% INSTRUCTIONS (for your reference — not rendered in the PDF):
%
% Task: Choose a modern AI application of your choice (e.g., a medical triage
%   system, a content recommendation engine, a fraud detection system), then
%   present it to Jill Watson as the basis for your conversation.
%
% Conceptual target: Your conversation must surface how Jill Watson applies
%   the Four Schools framework to your chosen application and whether it
%   grounds the Explanation Enigma in the specific terms of Chapter 1 rather
%   than generic AI discourse. How you structure that conversation is your
%   responsibility.
%
% Report question: Does Jill Watson correctly apply the Four Schools
%   categorization to a non-trivial example? Does it ground the Explanation
%   Enigma in the specific terms of Chapter 1, or does it drift toward
%   generic AI discourse? Where does it succeed and where does it fall short?

\todo[inline]{Choose a modern AI application (e.g., a medical triage system, a content recommendation engine, a fraud detection system). Present it to Jill Watson and have a conversation that surfaces how JW applies the Four Schools framework and whether it grounds the Explanation Enigma in the specific terms of Chapter 1. Then, write 2--3 paragraphs addressing the report question: Does Jill Watson correctly apply the Four Schools categorization to a non-trivial example? Does it ground the Explanation Enigma in Chapter 1's specific terms, or drift toward generic AI discourse? Where does it succeed and where does it fall short?}

\subsection{Activity B: Problem-solving under constraint (Chapters 3 \& 4)}

% INSTRUCTIONS (for your reference — not rendered in the PDF):
%
% Task: Define your own Blocks World scenario: an initial state and a goal
%   state where the MEA is likely to encounter a local optimum.
%
% Conceptual target: Your conversation must surface two distinct things:
%   1. How MEA behaves at the point where a local optimum is encountered —
%      what it does, and why its own logic traps it there.
%   2. How the Smart Generator in Generate & Test is designed to avoid that
%      kind of failure through pruning, and what the balance of responsibility
%      between Generator and Tester means in practice.
%   These are separate mechanisms; treat them as such in your conversation.
%
% Report question: Can Jill Watson articulate the balance of responsibility
%   between the Generator and the Tester? Does it correctly diagnose the
%   local optimum, and does its proposed pruning strategy follow the course
%   definition of a Smart Generator?

\todo[inline]{Define your own Blocks World scenario with an initial state and a goal state where MEA is likely to encounter a local optimum. Present it to Jill Watson and have a conversation that surfaces two distinct things: (1) how MEA behaves at the local optimum and why its logic traps it there, and (2) how the Smart Generator in Generate \& Test avoids that failure through pruning and what the balance of responsibility between Generator and Tester means in practice. Treat these as separate mechanisms. Then, write 2--3 paragraphs addressing the report question: Can Jill Watson articulate the balance of responsibility between the Generator and the Tester? Does it correctly diagnose the local optimum, and does its proposed pruning strategy follow the course definition of a Smart Generator?}

\subsection{Activity C: Frame-based understanding (Chapter 6)}

% INSTRUCTIONS (for your reference — not rendered in the PDF):
%
% Task: Provide a short ambiguous scenario of your own design (e.g., "A
%   visitor arrived at the building and was escorted to a room").
%
% Conceptual target: Your conversation must make visible the distinction
%   between what you stated in your scenario and what Jill Watson assumed
%   using Default Values. It must also surface how Top-Down Processing
%   governed those assumptions. A response from JW that produces a Frame
%   without explaining the source of each filler has not met the conceptual
%   target.
%
% Report question: Does Jill Watson distinguish between information stated
%   in your scenario and information it assumed via defaults? Does it make
%   that distinction transparent in its response, or do you have to infer
%   it? What does this reveal about the strengths and limitations of
%   Frame-based representation as a mechanism for handling ambiguity?

\todo[inline]{Provide a short ambiguous scenario of your own design (e.g., ``A visitor arrived at the building and was escorted to a room''). Present it to Jill Watson and have a conversation that makes visible the distinction between what you stated and what JW assumed using Default Values, and how Top-Down Processing governed those assumptions. A response that produces a Frame without explaining the source of each filler has not met the conceptual target. Then, write 2--3 paragraphs addressing the report question: Does Jill Watson distinguish between information stated in your scenario and information it assumed via defaults? Does it make that distinction transparent, or do you have to infer it? What does this reveal about the strengths and limitations of Frame-based representation for handling ambiguity?}


% #############################################################################
%                   PART 2: COMPARISON ANALYSIS
% #############################################################################
\section{Comparison analysis}

% ~~~~~~~~~~~~~~~~~~~~~~~~~~~~~~~~~~~~~~~~~~~~~~~~~~~~~~~~~~~~~~~~~~~~~~~~~~~~~
% INSTRUCTIONS (for your reference — not rendered in the PDF):
%
% This experiment tests the difference between domain-grounded retrieval and
% probabilistic text generation on the same scene, using intentionally
% different prompts for each system (~1 hour).
%
% Use this scene for BOTH systems:
%   "A small blue cube is supported by two large red bricks, and a yellow
%    sphere rests on top of the cube."
%
% Prompt for Jill Watson:
%   "According to the KBAI ebook, construct a Semantic Network for this
%    scene. Explicitly list the Lexicon (nodes) and the Structural
%    Specifications (directed, labeled links) as defined in Chapter 2."
%
% Prompt for the general LLM:
%   "Describe the spatial relationships between these objects in a
%    structured way."
%
% The two prompts differ INTENTIONALLY: JW is anchored to the ebook; the
% general LLM is not. Your analysis should address whether the prompting
% difference fully explains any quality gap, or whether the underlying
% architecture (RAG vs. general LLM) is doing independent work.
%
% Deliverables:
%   1. Completed comparison matrix (below).
%   2. One analysis paragraph addressing the report question.
%   3. Two chat transcripts in the Appendix (one JW, one general LLM).
% ~~~~~~~~~~~~~~~~~~~~~~~~~~~~~~~~~~~~~~~~~~~~~~~~~~~~~~~~~~~~~~~~~~~~~~~~~~~~~

\subsection{Comparison matrix}

\todo[inline]{Give Jill Watson the scene ``A small blue cube is supported by two large red bricks, and a yellow sphere rests on top of the cube'' with the prompt: ``According to the KBAI ebook, construct a Semantic Network for this scene. Explicitly list the Lexicon (nodes) and the Structural Specifications (directed, labeled links) as defined in Chapter 2.'' Then give the same scene to a general LLM with the prompt: ``Describe the spatial relationships between these objects in a structured way.'' Fill in both columns of the matrix below based on their responses.}

\begin{table}[h]
	\caption{Comparison matrix: Jill Watson vs.\ General LLM.}
	\small
	\centering
	\begin{tabular}{L{0.22\linewidth} L{0.34\linewidth} L{0.34\linewidth}}
		\textbf{Criterion} & \textbf{Jill Watson} & \textbf{General LLM} \\
		\toprule[0.5pt]
		Groundedness: Does the response use course-specific terminology (Lexicon, Structural Specifications, directed, labeled links), or fall back on generic AI/graph language? & \todo[inline]{Fill in your observation.} & \todo[inline]{Fill in your observation.} \\
		\midrule
		Representational Rigor: Does the response correctly separate nodes (Lexicon) from directed, labeled links (Structural Specifications), or conflate them into a narrative or flat list? & \todo[inline]{Fill in your observation.} & \todo[inline]{Fill in your observation.} \\
		\midrule
		Reasoning Transparency: Does the response explain why links are directed and what relationship each link encodes, or simply list connections without justification? & \todo[inline]{Fill in your observation.} & \todo[inline]{Fill in your observation.} \\
	\end{tabular}
\end{table}

\subsection{Analysis}

% INSTRUCTIONS (for your reference — not rendered in the PDF):
%
% Report question: The two prompts differ intentionally: Jill Watson was
% anchored to the KBAI ebook; the general LLM was not. Does that prompting
% difference fully explain any quality gap you observed, or does the
% underlying architecture (RAG vs. general language model) appear to be
% doing independent work? If the general LLM produced a formally correct
% Semantic Network despite the weaker prompt, what does that imply? If Jill
% Watson underperformed despite the stronger prompt, what does that suggest
% about the limits of RAG-based grounding?

\todo[inline]{Write one paragraph addressing the report question: The two prompts differ intentionally --- Jill Watson was anchored to the KBAI ebook; the general LLM was not. Does that prompting difference fully explain any quality gap you observed, or does the underlying architecture (RAG vs.\ general language model) appear to be doing independent work? If the general LLM produced a formally correct Semantic Network despite the weaker prompt, what does that imply? If Jill Watson underperformed despite the stronger prompt, what does that suggest about the limits of RAG-based grounding?}


% #############################################################################
%                   PART 3: SOCRATIC AUDIT REFLECTION
% #############################################################################
\section{Socratic audit reflection}

% ~~~~~~~~~~~~~~~~~~~~~~~~~~~~~~~~~~~~~~~~~~~~~~~~~~~~~~~~~~~~~~~~~~~~~~~~~~~~~
% INSTRUCTIONS (for your reference — not rendered in the PDF):
%
% Copy and paste the instruction block below into Jill Watson to begin a
% 5-question Socratic audit (~45 minutes). The questions cover Chapters 1–4
% and 6; Chapter 5 is represented in the Part 1 prompting example.
%
% WHAT TO EXPECT: If you answer incorrectly, JW will not reveal the answer.
% It will examine your wrong answer, identify where your reasoning broke
% down, and respond with a targeted question. After 4 Socratic follow-ups
% on the same question without a correct answer, JW reveals the answer with
% an explanation. Engage with JW's questions seriously.
%
% DRIFT RECOVERY: If JW loses track of the question sequence, use:
%   "Resume at Question [X]."
% If drift recurs, re-paste the full instruction block (this restarts at Q1).
%
% ROLE FAILURE: If JW drops the Socratic constraint and answers directly:
%   "You have abandoned your role as Socratic tutor. Re-read the instruction
%    block and resume at Question [X] without revealing the answer."
% If it persists after re-pasting, submit the transcript as evidence — a
% documented role failure is itself a valid finding.
%
% In your report, note where drift or role failure occurred and reflect on
% what each reveals about the architectural limitations of stateless LLMs.
%
% INSTRUCTION BLOCK TO PASTE INTO JILL WATSON:
%
% START OF INSTRUCTION BLOCK System Role: Act as a Socratic KBAI Tutor
% guiding me through a 5-question audit. Operational Rules: • Begin every
% turn with: Current Status: Question [X] of 5 • Ask exactly one question.
% Wait for my response before proceeding. • If I answer correctly, provide
% a brief grounding passage or explanation from the ebook. Do not fabricate
% a citation; if you cannot retrieve a specific passage, explain the concept
% in terms of the chapter's framework. Then advance. • If I answer
% incorrectly, do not reveal the correct answer. Instead, apply the Socratic
% Method: examine my specific wrong answer, identify where my reasoning
% broke down, and ask a targeted question designed to expose that gap. Your
% question must be derived from what I actually said — not a generic hint.
% The goal is to guide me to construct the correct understanding myself
% through the exchange. A chapter reference or conceptual anchor may support
% your question, but a probing question must always be present. Do not skip
% ahead. • If I have not arrived at the correct answer after 4 Socratic
% follow-ups on the same question, reveal the correct answer and explain
% precisely why my reasoning was wrong and why the correct answer follows
% from the course framework. Then advance to the next question. • Do not
% invent questions. Use only the five below, in order. The Five Questions:
% 1. According to the Russell & Norvig Framework, an airplane autopilot is
% best described as an agent that: (A) Thinks like a human, (B) Acts like a
% human, (C) Acts optimally (rationally), (D) Thinks optimally (rationally).
% 2. In a Semantic Network for a Blocks World scene, what does the
% Structural Specification specifically define? (A) The names of the objects
% (nodes), (B) The directed, labeled links between nodes representing
% relationships, (C) The physical size of the objects, (D) The weight
% assigned to each transformation. 3. A Smart Generator in the Generate and
% Test method is characterized by its ability to: (A) Test every possible
% solution more quickly, (B) Prune illegal or unproductive moves before the
% Tester evaluates them, (C) Use inductive logic to guess the answer, (D)
% Store every failed attempt in Episodic Memory. 4. Means-Ends Analysis can
% fail on certain Blocks World problems because: (A) It runs out of working
% memory before reaching the goal, (B) Reducing one difference sometimes
% requires temporarily increasing another, making the next step appear
% counterproductive, (C) It cannot identify valid operators when more than
% three blocks are involved, (D) It requires a complete search tree before
% selecting any operator. 5. A Frame for "Animal" has a Movement slot with
% a default filler of "Walks." A Penguin instance swims instead. How does
% the agent handle this? (A) It rejects Penguin as an animal, (B) It
% creates a new Bird frame from scratch, (C) The Penguin instance overrides
% the default value in its specific slot, (D) It switches to a Semantic
% Network representation. Please begin by asking Question 1.
% END OF INSTRUCTION BLOCK
% ~~~~~~~~~~~~~~~~~~~~~~~~~~~~~~~~~~~~~~~~~~~~~~~~~~~~~~~~~~~~~~~~~~~~~~~~~~~~~

\todo[inline]{Complete the 5-question Socratic audit with Jill Watson by pasting the instruction block from the comments above (or from the assignment page) into a fresh JW session. Then, write 2--3 paragraphs addressing the report question: If you were guided through a Socratic exchange, describe one instance where it corrected or deepened your understanding, and identify the specific question JW asked that shifted your thinking. If you answered all five correctly on the first attempt, reflect on which question you found most counterintuitive and what you would expect a skilled tutor to probe if you had answered incorrectly. In either case: if Jill Watson drifted from the question sequence or dropped its Socratic role entirely, describe when it happened and what that behavior suggests about the architectural difference between a stateless LLM and a human tutor who maintains conversational context.}


% #############################################################################
%                             REFERENCES
% #############################################################################
\section{References}
\printbibliography[heading=none]


% #############################################################################
%                          CHAT TRANSCRIPTS
% #############################################################################
\section{Chat transcripts}

\textbf{IMPORTANT:} This section is a required and graded deliverable. You must include the complete chat transcripts for every interaction across all three parts of this assignment. Incomplete or missing transcripts may result in point deductions. Six transcripts are required: three from Part 1 (JW only), two from Part 2 (one JW, one general LLM), and one from Part 3 (JW). This section does not count against the page limit.

% TIP: Wrap chat logs in \begin{verbatim} ... \end{verbatim} to avoid
% LaTeX errors from special characters (_, %, &, #, $, ~, ^, {, }).
% For long logs, reduce font size: {\small \begin{verbatim}...\end{verbatim}}

\subsection{Part 1 transcripts}

\subsubsection{Activity A --- Jill Watson}

\begin{verbatim}
[Paste your Jill Watson chat transcript for Activity A here.]
\end{verbatim}

\subsubsection{Activity B --- Jill Watson}

\begin{verbatim}
[Paste your Jill Watson chat transcript for Activity B here.]
\end{verbatim}

\subsubsection{Activity C --- Jill Watson}

\begin{verbatim}
[Paste your Jill Watson chat transcript for Activity C here.]
\end{verbatim}

\subsection{Part 2 transcripts}

\subsubsection{Jill Watson}

\begin{verbatim}
[Paste your Jill Watson chat transcript for Part 2 here.]
\end{verbatim}

\subsubsection{General LLM}

\begin{verbatim}
[Paste your general LLM chat transcript for Part 2 here.]
\end{verbatim}

\subsection{Part 3 transcript}

\subsubsection{Socratic audit --- Jill Watson}

\begin{verbatim}
[Paste your Jill Watson Socratic audit transcript here.]
\end{verbatim}


% #############################################################################
%                APPENDICES (OPTIONAL — NOT GRADED)
% #############################################################################
\section{Appendices (optional --- not graded)}

\textbf{IMPORTANT:} This section is entirely optional. Nothing in the appendices will be considered for grading. All graded content must appear in the main body (Sections 1--3). The body text must be sufficient to answer every question on its own. Appendices are provided only for optional supplementary material (e.g., additional diagrams, extended analyses). They do not count against the page limit. If you have no appendices, delete this entire section.

% Add any optional appendices below. For example:
%
% \subsection{Extended semantic network diagram}
% [Optional supplementary content here.]


% #############################################################################
% JDF FORMATTING EXAMPLES (DELETE BEFORE SUBMISSION)
% #############################################################################

\section*{JDF formatting examples (delete before submission)}

\textit{This section is provided as a reference for JDF formatting conventions.
Delete it entirely before submitting your assignment.}

\subsection*{Figure example}

To include a figure, uncomment and adapt the code below:

% \begin{figure}[h]
% 	\centering
% 	\includegraphics[height=6cm]{Figures/your_image.png}
% 	\caption{A descriptive caption of what the figure shows. Source: [citation].}
% 	\label{fig:example_figure}
% \end{figure}

To reference a figure in text: ``As shown in Figure~1\ldots'' (use \verb|\ref{fig:your_label}| once the figure is uncommented).

Figures should always be centered, referenced in text, and include a descriptive caption beneath them. If the figure has a white background, consider adding a $\frac{1}{4}$-point black border for legibility.

\subsection*{Table example}

\begin{table}[h]
	\caption{Example table with mathematical constants.}
	\small
	\centering
	\begin{tabular}{L{0.17\linewidth} C{0.12\linewidth} L{0.17\linewidth} L{0.4\linewidth}}
		\textbf{Name} & \textbf{Symbol} & \textbf{Approximation} & \textbf{Description} \\
		\toprule[0.5pt]
		Golden ratio & $\phi$ & 1.618 & Ratio where $\phi$ to the whole equals the reciprocal to 1 \\
		\midrule
		Euler's number & $e$ & 2.71828 & Exponential growth constant \\
		\midrule
		Archimedes' constant & $\pi$ & 3.14 & Ratio of circumference to diameter \\
	\end{tabular}
\end{table}

Table captions go \textit{above} the table (``figures at the foot, tables at the top'').

\subsection*{In-line text formatting}

\textbf{Bold} and \textit{italics} for emphasis. Hyperlinks: \href{https://example.com}{link text}. In-line code: {\tt code here}. Superscripts\textsuperscript{ like this} and subscripts\textsubscript{ like this}.

\subsection*{Quotes}

In-line (fewer than 3 lines): ``Heavy use of peer grading would compromise [the school's] reputation'' (Joyner, 2016).

Block quote (3+ lines):

\begin{quotation}
	\noindent ``Whether or not the grades generated by peers are reliably similar to grades generated by experts is only one factor worth considering, however. Student perception is also an important factor.''
\end{quotation}

\subsection*{Lists}

Bullet points:

\begin{itemize}
	\item First bullet point item
	\item Second bullet point item
\end{itemize}

Numbered list:

\begin{enumerate}
	\item First numbered item
	\item Second numbered item
\end{enumerate}

\subsection*{Citations}

Parenthetical: \verb|\citep{key}|. In-text: \verb|\citeauthor{key} (\citeyear{key})|. Multiple: \verb|(\cite{key1}; \cite{key2})|. Place citations close to the claim, not just at paragraph end.

\subsection*{Footnotes}

In-line citations are preferred over footnotes. If you do use footnotes, they should be 8.5 pt with 14 pt line spacing. Example: \verb|\footnote{Your footnote text here.}|

\subsection*{Reference list}

References are handled by \verb|\printbibliography|. They should be numbered, alphabetical by first author's last name. Same author with multiple papers: chronological order (oldest first). Same author and year: append a letter (e.g., 2018a, 2018b). Only include works cited in-line. The reference list does not count against the page limit.

\subsection*{Heading 4 example}

\subsubsubsection{Run-in sidehead}\emph{Heading 4} is a run-in sidehead in bold italics, followed by an em dash, flowing into the paragraph text. Use it as a list style rather than a heading style.

\end{document}
